\documentclass{article}

\usepackage{amsmath}
%\usepackage{amsfonts}
\usepackage{amsthm}
%\usepackage{amssymb}
%\usepackage{mathrsfs}
%\usepackage{fullpage}
%\usepackage{mathptmx}
%\usepackage[varg]{txfonts}
\usepackage{color}
\usepackage[charter]{mathdesign}
\usepackage[pdftex]{graphicx}
%\usepackage{float}
%\usepackage{hyperref}
%\usepackage[modulo, displaymath, mathlines]{lineno}
%\usepackage{setspace}
%\usepackage[titletoc,toc,title]{appendix}
\usepackage{natbib}
\usepackage[makeroom]{cancel}
\usepackage[only,llbracket,rrbracket]{stmaryrd}

%\linenumbers
%\doublespacing

\theoremstyle{definition}
\newtheorem*{defn}{Definition}
\newtheorem*{exm}{Example}

\theoremstyle{plain}
\newtheorem*{thm}{Theorem}
\newtheorem*{lem}{Lemma}
\newtheorem*{prop}{Proposition}
\newtheorem*{cor}{Corollary}

\newcommand{\argmin}{\text{argmin}}
\newcommand{\ud}{\hspace{2pt}\mathrm{d}}
\newcommand{\bs}{\boldsymbol}
\newcommand{\PP}{\mathsf{P}}
\let\divsymb=\div % rename builtin command \div to \divsymb
\renewcommand{\div}[1]{\operatorname{div} #1} % for divergence
\newcommand{\Id}[1]{\operatorname{Id} #1}

\title{Of icy streams and fiery torrents}
\author{Daniel R. Shapero, Robert J. Baraldi}
\date{}

\begin{document}

\maketitle

Simulating phase change is a challenging free boundary problem that arises in multiple disciplines in the earth sciences.
The goal of this paper is to describe a new approach to numerical simulation of phase change based on the principle of convex duality.
We will then apply this method to simulating the flow of glaciers and lava.

The physical law we focus on is global conservation of enthalpy $E$.
The enthalpy is related to the temperature $T$ of the medium.
There are two mechanisms of transport: advection by some bulk velocity $u$ and diffusion.
We assume a general form for volumetric heat sources $Q$ and boundary fluxes $f$ because these will depend on the application.
For both glacier and lava flow, there are internal heat sources due to viscous dissipation of mechanical energy.
Lava flows lose heat to the atmosphere at the upper surface by radiation, but can eventually cool to the point of forming an insulating rocky crust at the surface.
Glaciers gain heat at their upper surface by insolation and the resulting meltwater then runs off the surface.

The first principle we work from is energy conservation.
Let $\omega$ be a representative volume element.
We will want to consider the possiblity that $\omega$ is changing in time.
Let $v$ be the velocity of the boundary of $\omega$ and $n$ the unit outward-pointing normal vector.
Then the usual transport theorem implies that the rate of change of total enthalpy within $\omega$ is a sum of the flux across the boundary and sources:
\begin{equation}
    \frac{\mathrm d}{\mathrm dt}\int_\omega E\ud x + \int_{\partial\omega}\left\{E(u - v) - k\nabla T\right\}\cdot n\ud s = \int_\omega Q\ud x + \int_{\partial\omega\cap\Gamma}f\cdot n\ud s.
\end{equation}
Again, we emphasize that the factor of $v$ accounts for the fact that the representative volume element $\omega$ can depend on time.

The main challenge comes from the relation between temperature and enthalpy.
For a material that does not change phase, the enthalpy is equal to $\rho c T$ where $\rho$ is the material density and $c$ is the specific heat capacity at constant pressure.
For a material that can change phase, say from solid to liquid, supplying more heat at the melting point does not raise the temperature but instead supplies the latent heat of transformation.
The effect is to make the enthalpy a multi-valued function of temperature:
\begin{equation}
    E(T) = \begin{cases}\rho_s c_s T & T < T_m \\ \rho_l c_l T + \rho_s L & T > T_m \end{cases}
\end{equation}
where $L$ is the latent heat of melting.
We can likewise invert the relationship to express the temperature as a single-valued function of enthalpy:
\begin{equation}
    T(E) = \begin{cases}E / \rho_sc_s & E \le \rho_sc_sT_m \\ T_m & \rho_sc_sT_m < E \le \rho_s(c_sT + L) \\ T_m + (E - \rho_sc_sT_m) / \rho_lc_l\end{cases}
\end{equation}
The form of the temperature-enthalpy relationship means that the problem is not only nonlinear, but the nonlinearity is discontinuous.

There are a number of by now orthodox methods for solving the Stefan problem.
First, we can eliminate the temperature from the problem.
The $\mathrm dE/\mathrm dt$ term in the conservation law is then easy to handle, but the diffusion term $\nabla\cdot k\nabla T(E)$ becomes difficult because $T(E)$ is only piecewise-differentiable.
Using a variational form allows us to paper over that difficulty at first.
But in practice the lack of differentiability makes nonlinear solvers converge slowly or not at all.
The remedy is then to regularize the temperature-enthalpy relationship.
That approach then leads us to an alternative based on phase-field methods.
Here an additional field $\phi$ is introduced representing which material phase is present at a given point in space.
The phase field becomes a coefficient in the heat equation that can be used to interpolate the density, specific heat, thermal conductivity, and latent heat between the two material phases.
We then devise a balance law for $\phi$ that includes a small parameter $\epsilon$ which controls how much the solid-liquid phase interface is smeared out.
In the limit as $\epsilon \to 0$, we should recover the original Stefan problem.
The challenge with phase field methods is that the differential equation for $\phi$ is strongly nonlinear and becomes even more stiff than the diffusion equation.

Here we propose to instead include both $T$ and $E$ as unknowns in the problem.
In so doing, we will recover in a modified form the gradient flow structure of the heat equation without phase change.

A differential equation for a field $z$ is a gradient flow when there exists some \emph{free energy functional} $F$ such that
\begin{equation}
    \frac{\mathrm dz}{\mathrm dt} = -\nabla F(z).
\end{equation}
A simple argument using the chain rule shows that the free energy decreases along trajectories.
With some mild regularity assumptions, the free energy is a global Lyapunov functional, which implies that the system is stable.

This definition of a gradient flow in the purest sense is overly strict.
For example, if we discretize a gradient flow in space with finite elements, we instead end up with a finite-dimensional system of the form
\begin{equation}
    M\frac{\mathrm dz}{\mathrm dt} = -\nabla F(z)
\end{equation}
where $M$ is the finite element mass matrix.
Strictly speak, we can multiply by $M^{-1}$ and define a new, effective free energy.
But it would be better if we could avoid dirty hacks like that.
There is a more general class of problems, however, called \emph{generalized} gradient flows.
A differential equation is a generalized gradient flow if there exists a convex functional $G$ such that
\begin{equation}
    \frac{\mathrm d}{\mathrm dt}\nabla G(z) = -\nabla F(z).
\end{equation}
We recover the class of ordinary gradient flows when we take $G(z) = \|z\|^2/2$ and the class of gradient flows with a symmetric, positive-definite matrix by taking $G(z) = \|z\|_M^2 / 2$.
Provided that $G$ is strongly convex, i.e. the second derivative $\nabla^2G$ exists and is strictly positive-definite, then the same argument as above shows that $F$ is decreasing along trajectories.

Suppose that we discretize a generalized gradient flow with the backward Euler method.
Then after some rearranging the resulting system of equations to solve for the next state $z_{n + 1}$ from the current state $z_n$ is
\begin{equation}
    \nabla G(z_{n + 1}) - \nabla G(z_n) + \delta t\cdot \nabla F(z_{n + 1}) = 0.
\end{equation}
But this nonlinear system of equations is also an optimization problem:
\begin{equation}
    z_{n + 1} = \text{minimize}_z\left\{G(z) - \langle \nabla G(z_n), z\rangle + \delta t\cdot F(z)\right\}.
\end{equation}
The objective is a kind of generalized proximal operator, which motivates the connection of the backward Euler discretization to proximal methods in optimization.

The heat equation without advection is a generalized gradient flow with
\begin{equation}
    G(T) = \frac{1}{2}\int_\Omega \rho cT^2\ud x, \quad F(T) = \int_\Omega\left(\frac{1}{2}k|\nabla T|^2 - Q\cdot T\right)\ud x.
\end{equation}
The heat equation with phase change is also a generalized gradient flow, with the caveat that the functional $G$ is no longer smooth because the temperature-enthalpy relationship is non-differentiable.
At a purely mechanical level, trying to determine the minimization or gradient flow form of a differential equation amounts to evaluating antiderivatives.
Since $E(T)$ is piecewise linear, we expect that its antiderivative is piecewise quadratric, with a discontinuity in slope at $T_m$:
\begin{equation}
    \mathscr G(T) = \begin{cases}\frac{1}{2}\rho_sc_s(T - T_m)^2 & T \le T_m \\ \frac{1}{2}\rho_lc_l(T - T_m)^2 + \rho_sL(T - T_m) & T > T_m\end{cases}
\end{equation}
We could then try to devise a numerical scheme for the Stefan problem by discretizing in time with the backward Euler method.
But if we do so we are right back where we started: we have a large system of nonsmooth nonlinear equations.
We know more about its solution in the form of the gradient flow structure but we are no closer to a method to approximating this solution numerically.

The crux of this paper is that we \textbf{can work with \emph{both} the temperature and enthalpy as unknowns}.
The highest level at which we can view this approach is through convex duality.
We assumed that the functional $G$ in the generalized gradient flow above is convex, so it has a convex conjugate functional $G^*$.
The relationship between $G$ and $G^*$ is that their derivatives are inverses of each other:
\begin{equation}
    w = \nabla G(z) \Leftrightarrow z = \nabla G^*(w).
\end{equation}
We can then rewrite the generalized gradient flow as a \emph{system} of equations:
\begin{align}
    \frac{\mathrm dw}{\mathrm dt} & = -\nabla F(z) \\
    0 & = \nabla G^*(w) - z\nonumber
\end{align}
To our knowledge, there is no standard terminology for this problem, so we refer to it as the \emph{dual} gradient flow.
But first we describe the discretization.
If we again use the backward Euler method for the differential-algebraic equation above, we arrive at the system
\begin{align}
    w_{n + 1} - w_n + \delta t\cdot \nabla F(z_{n + 1}) & = 0 \\
    z_{n + 1} - \nabla G^*(w_{n + 1}) & = 0.
\end{align}
Just like how the backward Euler discretization of a generalized gradient flow is equivalent to a minimization problem, the backward Euler discretization of a dual gradient flow is a saddle point problem:
\begin{equation}
    z_{n + 1},\; w_{n + 1} = \text{maximize}_w\,\text{minimize}_z\;\left\{\langle w - w_n, z\rangle - G^*(w) + \delta t\cdot F(z)\right\}.
\end{equation}
Saddle-point problems are more challenging to solve numerically than pure minimization problems.
But depending on the nature of the nonlinearity, $G^*$ can be so much more convenient to work with than $G$ that the extra cost of the saddle-point problem is worth it.
As we argue below, this is the case for the Stefan problem.

We can compute the conjugate of the enthalpy-temperature relationship by hand:
\begin{equation}
    \mathscr G^*(E) = \begin{cases}
        \frac{1}{2\rho_sc_s}(E - \rho_sc_sT_m)^2 & E \le \rho_sc_sT_m \\
        0 & \rho_sc_sT_m \le E \le \rho_s(c_sT_m + L) \\
        \frac{1}{2\rho_lc_l}(E - \rho_s(c_sT_m + L))^2 & E \ge \rho_s(c_sT_m + L)
    \end{cases}
\end{equation}
In the gradient flow formulation of the Stefan problem above, we worked with the multi-valued mapping from temperature to enthalpy.
By passing to the convex conjugate, we get to instead work with the single-valued mapping from enthalpy to temperature.

There is some similarity between the Stefan problem and other degenerate parabolic equations like the porous medium equation \citep{vazquez2006porous}.


\pagebreak

\bibliographystyle{plainnat}
\bibliography{stefan-problem.bib}

\end{document}
